\scanpage{70}
Or choisissons un $\zeta\in A^1(X)$ de polarisation, Alors on a un diagramme comm.
\diagram{\begin{tikzcd}[ampersand replacement=\&,column sep=large,row sep=large]
H^{2n-1}(X,Q_l)(n-1)\arrow[r,"{\substack{\psi_X\\\sim}}"] \& H^1(B_X,Q_l)\\
H^1(X,Q_l)\arrow[u,"{\lambda\quad\unsure{\wr|2}}","{\text{cup }\zeta^{n-1}}"'] \& H^1(A_X,Q_l)\arrow[l,"\varphi_X^{\deleted{\unsure{(1)}}}","\sim"']\arrow[l,phantom,"{\deleted{\rightarrow}}"{description,pos=.08}]\arrow[u,swap,"{\wr\quad\mu}"]
\end{tikzcd}}
où toutes les flèches écrites \deleted{\unsure{sont}} \deleted{\unsure{et}} définies par des classes de corr. alg. (ind. de $l$). \inserted{($\mu$ étant défini comme composé \inserted{$\psi_X\lambda\varphi_X$}) et sont des isom. en vertu de \ldots{} et de 2.6.} Donc $\psi_X^{-1}=\lambda\varphi\mu^{-1}$, et tout revient à voir que $\mu^{-1}$ est algébrique, ce qui provient du fait que $\mu$ est défini par un $Q$-\unsure{isom|hom|isogénie} $B_X\unsure{\mathbin{-}|\to}A_X$ \inserted{(cf}, dont on prend simplement l'inverse\ldots{}

Reste à traiter le cas de $\varphi_Y^{-1}$ (où on peut remplacer $X$ par $Y$). On utilise le $\unsure{\overline{\mathrm m}|=}$ diagramme, qui donne
\[
\varphi_X^{-1}=\mu^{-1}\psi_X\lambda
\]
donc on est ramené à inverser $\mu$ \unsure{``algébriquement''|``algébriquemt''|``algébriques''}, or cela est possible grâce au fait que $\mu$ est associé à une isogénie de V.A.
% archive sheet number: 69

